\section*{ID8137: What I Attempted to Add: E≈ -V×B}
\paragraph{Metadata.}
PiID: 7817113864558 | RTEID: RTE:P31597.7780:T1.7979 | UUID: 36ed43b4-9015-552e-8c31-63f06e395746

\paragraph{Canonical Statement.}
- **(B) Gyro comoving observer \textbackslash{}(u\textasciicircum{}\textbackslash{}mu\textbackslash{})** (gyro rest frame): same covariant \textbackslash{}(T\textasciicircum{}\{\textbackslash{}mu\textbackslash{}nu\}\textbackslash{}) but contracted with \textbackslash{}(u\textasciicircum{}\textbackslash{}mu\textbackslash{}). Motion induces electric fields from the dipole’s magnetic field (\textbackslash{}(\textbackslash{}mathbf\{E\}\textbackslash{}approx -\textbackslash{}mathbf\{v\}\textbackslash{}times\textbackslash{}mathbf\{B\}\textbackslash{})), boosting the local energy density and changing Poynting flux directions; some off-diagonals can reduce in a comoving tetrad while anisotropic mechanical stresses remain.

\paragraph{Canonical Equation.}
\[
\mathbf{E}\approx -\mathbf{v}\times\mathbf{B}
\]

\paragraph{Dictionary.}
Relation: E≈ -v×B

\paragraph{Encyclopedia.}
What I Attempted to Add: E≈ -V×B is a scaling / order-of-magnitude relation, written E≈ -v×B. It expresses a scaling or proportionality, capturing how the quantity grows with its parameters. It is built from E, v, B. Within the Trawin Zero Point registry it serves as one element of the framework's mathematical narrative.
